% Working method description bound to protocol-v2; not a final evidence claim.
\section{Problem Formulation}
We study continuous price prediction for advertised used vehicles. A source
record contains covariates $x_i$, their original missingness pattern $m_i^0$,
and a training outcome $y_i=\log(1+p_i)$, where $p_i$ is the recorded advertised
price. This outcome is neither a transaction price nor an independently
measured intrinsic value. A reported listing $z_i(H_i)$ retains the original
missing entries and additionally withholds a set $H_i$ of originally recorded
fields. In the primary protocol, the optional fields are mileage,
transmission, fuel type, and body condition. The feasible report set is
\begin{equation}
\mathcal H_i=\{H\subseteq\mathcal A\setminus m_i^0:\ |H|\leq h\},
\qquad h=2,
\label{eq:feasible_reports}
\end{equation}
In Eq.~\eqref{eq:feasible_reports}, $\mathcal A$ contains the four optional
fields. The eleven candidate
masks include the empty mask, four singletons and six pairs; masks that
attempt to hide an originally absent value are invalid for that record.
Identifiers and other recorded covariates cannot be altered. This threat
model concerns withholding, rather than fabrication or changes to a vehicle.

The honest report is $z_i(\varnothing)$, which can still contain natural
missingness. The platform observes the reported covariates and their visible
missingness; it does not observe $H_i$, $m_i^0$, unavailable source values or
evaluation outcomes. Thus, the same reported listing must receive the same
pre-query scores whether a missing mileage entry is natural or concealed.
For example, a recorded mileage of 140,000 km that is withheld and an
originally unavailable mileage entry can produce identical public inputs.
The mileage number is illustrative. The distinction becomes visible only
through the stipulated query response, not through a secret eligibility flag.

A query requests one visibly missing field. Let $T_a(z_i(H_i);x_i)$ denote
the replay transition for field $a$. If its source value exists, the query
restores that recorded value; otherwise the transition is the identity.
Both outcomes consume one query slot. This is an oracle replay of available
records, not physical inspection or authentication of seller claims.
Each batch of $n$ listings has capacity $B=\lfloor\rho n\rfloor$, and at most
one field may be queried per listing. There is no borrowing of unused slots
between batches.

For valuation function $f_\theta$ and allocation rule $\pi_\phi$, the
stipulated seller payoff is
\begin{equation}
u_i(H_i,H_{-i})=
f_\theta\!\left(T_{a_i(H_i,H_{-i})}(z_i(H_i);x_i)\right)
-\kappa|H_i|,
\label{eq:seller_payoff}
\end{equation}
Equation~\eqref{eq:seller_payoff} uses the identity transition when no query
is assigned. Sellers know the
deterministic scoring and batch-allocation rule. Their payoff is a simulated
log-quote objective with concealment cost $\kappa$, rather than an inferred
behavioral utility. Capacity couples the sellers because changing one report
can change which listings are queried. Strategic classification and
performative prediction already establish the need to account for responses
to deployed rules~\cite{hardt2016strategic,perdomo2020performative}; we do not
claim that auditing or feedback itself is new.

\section{Method}
SAV-Val combines mask-aware valuation, a suppression-output query policy,
and policy-aware residual training. Its candidate contribution is the
empirical value of this particular combination under shared capacity and
ambiguous missingness. An advantage over the controls, disclosure incentive compatibility,
equilibrium existence and convergence are not consequences of the design;
they require separate evidence or additional assumptions.

\subsection{Mask-aware valuation}
A CatBoost backbone~\cite{prokhorenkova2018catboost} is fitted on the feasible
reported states of each estimator-training record. Every state retains the
record's outcome, and each state has weight $1/|\mathcal H_i|$, so records with
more feasible masks do not receive more total training weight. A two-layer
ReLU residual network with 64 units per hidden layer adjusts the frozen
backbone:
\begin{equation}
f_\theta(z)=b(z)+r_\psi(e(z)).
\label{eq:residual_valuation}
\end{equation}
The encoder $e$ in Eq.~\eqref{eq:residual_valuation} is fitted using training
covariates only. Numeric features use
median imputation, missing indicators and scaling; categorical features use
one-hot encoding with an explicit missing token. Mileage and supported
range/capacity variables are log-transformed before encoding. Outcomes,
identifiers and declared direct reveal aliases are excluded. Missingness-aware
prediction has established theoretical and architectural
antecedents~\cite{lemorvan2020neumiss,lemorvan2021imputation}; this residual
construction is not claimed to reproduce their guarantees.

\subsection{Suppression-target construction}
For an offline training record and report $z=z_i(H)$, define the positive
quote gap relative to its honest report as
\begin{equation*}
d_i(z)=[f_\theta(z)-f_\theta(z_i(\varnothing))]_+.
\end{equation*}
The suppression target for a visibly missing field is
\begin{equation}
s_{ia}(z)=d_i(z)-d_i(T_a(z;x_i)).
\label{eq:suppression_target}
\end{equation}
The honest-reference and post-query quotes in
Eq.~\eqref{eq:suppression_target} form training targets only.
They are absent from runtime head inputs. A failed replay has zero
suppression target. A successful replay can also have a nonpositive target;
revealing a value need not lower a quote.

For comparison, an ordinary error-benefit target is
\begin{equation}
v_{ia}(z)=(f_\theta(z)-y_i)^2
-(f_\theta(T_a(z;x_i))-y_i)^2.
\label{eq:benefit_target}
\end{equation}
The ordinary error reduction in Eq.~\eqref{eq:benefit_target} supplies a
matched target to the implemented Extra-Trees head~\cite{geurts2006extratrees},
which jointly predicts
three outputs: suppression $s_{ia}$, error benefit $v_{ia}$, and before-query
squared error. Its inputs are the encoded public report, its current quote
and a one-hot candidate field. The primary policy ranks the suppression
output, while the ordinary-benefit counterpart ranks the benefit output.
Therefore, this is a comparison of policy outputs and their induced training
feedback, not a claim that the primary head is trained exclusively on
suppression labels. Query candidates include naturally missing fields; the
head cannot receive the simulator's deliberate-concealment mask.

Active acquisition has previously been learned through reinforcement
learning, generative surrogates, mutual-information objectives, and
nongreedy oracles~\cite{janisch2019costly,li2021gsm,covert2023information,valancius2024aco}.
Our one-query ranking policy is a local shared-capacity design and is not an
official implementation of these sequential acquisition methods.

\subsection{Budget allocation}
For each listing, let $g_{ia}$ be its predicted score for a visibly missing
field. Choose the highest-scoring field, then rank the listings by those
scores. Query at most the first $B$ listings with strictly positive finite
scores. Ties are broken by fixed within-batch order. This procedure solves
the additive surrogate allocation problem in Eq.~\eqref{eq:surrogate_allocation}:
\begin{align}
\max_{q}\quad &\sum_{i,a}q_{ia}g_{ia},\nonumber\\
\text{subject to}\quad &q_{ia}\in\{0,1\},\quad
\sum_a q_{ia}\leq1,\quad\sum_{i,a}q_{ia}\leq B,
\label{eq:surrogate_allocation}
\end{align}
with $q_{ia}=0$ for publicly observed fields. To see this, replace any selected
field by the best field of that listing without decreasing the objective,
then exchange any selected lower-score listing for an unselected
higher-score listing. The result is the top-$B$ positive ranking.
This conditional surrogate optimality does not imply optimal realized price
error, optimal seller incentives, or global cross-batch allocation.

\subsection{Policy-aware feedback}
The residual is trained with honest-report error, mask-averaged error,
attacked-report error, a positive-gap hinge, and an honest-backbone anchor.
Writing $\widehat z_i$ for the report after the simulated allocation,
the implemented objective is Eq.~\eqref{eq:training_loss}:
\begin{align}
\mathcal L={}&\mathbb E_i[(f_\theta(z_i(\varnothing))-y_i)^2]
+0.25\mathcal L_{\rm mask}\nonumber\\
&+0.5\mathbb E_i[(f_\theta(\widehat z_i)-y_i)^2]\nonumber\\
&+\beta\mathbb E_i\left[\left(
f_\theta(\widehat z_i)-\operatorname{sg}(f_\theta(z_i(\varnothing)))
-\tau\right)_+\right]\nonumber\\
&+0.25\mathbb E_i[(f_\theta(z_i(\varnothing))-b(z_i(\varnothing)))^2],
\label{eq:training_loss}
\end{align}
where $\mathcal L_{\rm mask}$ is the per-record mean squared error over its
feasible masks, $\operatorname{sg}$ stops gradients, $\beta=0.5$, and
$\tau=0.02$ log-price units. The anchor and honest error discourage moving the
reference quote merely to reduce measured inflation. They do not eliminate
that possibility, so signed bias and honest-price error remain necessary
diagnostics.

Within each fitting loop, estimator and feedback-guidance records are
disjoint. The backbone fits estimator records, while the surrogate head fits
guidance records. Two feedback rounds each train the residual for six epochs
and refresh guidance targets. Query scores are refreshed once per epoch and
held fixed within it. Selected report indices and allocation are discrete;
gradients pass through their selected price predictions. This is a surrogate
alternating procedure, not differentiation through an exact bilevel game.
Training uses two seller-response sweeps and one start. Evaluation uses a
stronger independently checked response search.

\subsection{Nested target generation}
Three balanced outer folds are assigned by outcome-independent hashes of
listing identifiers. Each outer teacher fits the \emph{entire} backbone,
surrogate-head and residual-feedback loop on its outer-training records,
using a further 70/30 estimator/guidance split. It predicts every candidate
report on its outer-held-out records without receiving their outcomes.
The collected outer quotes, original training outcomes and replay transitions
then form targets for deployment heads.

The deployment estimator is fitted separately using the same 70/30 inner
split on the full allowed training cohort. The deployment head learns from
the collected outer targets, using public-state encodings from that
deployment estimator. Consequently, outer-held-out outcomes cannot influence
the corresponding teacher's fitted feedback loop. Nevertheless, outer
teachers and the deployment estimator use different fitting sample sizes and
encoders. The procedure gives an honest target-generation boundary, rather
than exact equality between teacher and deployment prediction distributions.

\subsection{Seller-response search and deviation audit}
Evaluation performs coordinate best responses with four starts: honest
reports, individual pre-query quote maximizers, and two seeded feasible
report profiles. At each update, the seller's candidate is ranked against
fixed peer reports under the actual capacity and tie rule. Ties retain the
current report; repeated profiles are detected. The search uses at most
20 sweeps per start and retains the profile with largest aggregate simulated
seller payoff. This selection does not prefer convergence over attack value.

An independent checker then enumerates every feasible unilateral deviation
and reruns the allocator. It verifies the search's remaining deviation regret,
\begin{equation*}
r_i=\max_{H\in\mathcal H_i}u_i(H,H_{-i})-u_i(H_i,H_{-i}),
\end{equation*}
and records the largest unilateral quote at those fixed peer reports.
The checker can be exact for this finite unilateral comparison even when the
joint response search has not converged. Its quote envelope is neither a
jointly attainable batch outcome nor a bound over all peer profiles. Both
remaining regret and nonconverged batches must accompany performance claims.
